%% For double-blind review submission
\documentclass[acmlarge,fleqn,review]{acmart}\settopmatter{printfolios=true}
%% For single-blind review submission
%\documentclass[acmlarge,review]{acmart}\settopmatter{printfolios=true}
%% For final camera-ready submission
%\documentclass[acmlarge]{acmart}\settopmatter{}

%% Note: Authors migrating a paper from PACMPL format to traditional
%% SIGPLAN proceedings format should change 'acmlarge' to
%% 'sigplan,10pt'.

%include polycode.fmt

\newcommand{\keyword}[1]{\mathbf{#1}}
\newcommand{\identifier}[1]{\mathit{#1}}
\makeatletter
\newcommand{\shorteq}{%
  \settowidth{\@@tempdima}{-}%
  \resizebox{\@@tempdima}{\height}{=}%
}
\makeatother

%format of = "\keyword{of}"
%format = = "{=}"
%format -> = "{\to}"
%format => = "{\Rightarrow}"
%format | = "{|}"
%format == = "{\shorteq\kern.5pt\shorteq}"
%format /= = "{\neq}"
%format && = "{\wedge}"
%format || = "{\vee}"
%format : = "{:}"
%format `elem` = "{\in}"
%format . = "{\circ}"
%format ~ = " "

%format (all(x)) = "\forall" x "."
%format (some(x)) = "\exists" x "."
%format empty = "\emptyset"
%format <~ = "{\hookleftarrow}"
%format (LUB(x)) = "\bigsqcup_{" x " \in \mathbb N}\kern-\smallskipamount"

%format fail = "\keyword{fail}"
%format skip = "\keyword{skip}"
%format replace = "\keyword{replace}"
%format rearrS = "\keyword{rearrS}"
%format rearrV = "\keyword{rearrV}"
%format case = "\keyword{case}"
%format normal = "\keyword{normal}"
%format exit = "\keyword{exit}"
%format exit' = "\identifier{exit}"
%format adaptive = "\keyword{adaptive}"

\definecolor{putback}{RGB}{60,116,219}
\definecolor{range}{RGB}{0,160,0}

\newcommand{\assert}[1]{{\color{putback}\{}\,#1\,{\color{putback}\}}}
\newcommand{\varassert}[2]{{\color{putback}\{\,#1\kern-\smallskipamount:}\ \ #2\,{\color{putback}\}}}
\newcommand{\assertrange}[1]{{\color{range}\{\!\{}\,#1\,{\color{range}\}\!\}}}
\newcommand{\varassertrange}[2]{{\color{range}\{\!\{\,#1\kern-\smallskipamount:}\ \ #2\,{\color{range}\}\!\}}}

%format (ENV(x)) = "\overrightarrow{" x "}"
%format (COM(x)) = "\langle\," x "\,\rangle"
%format (ASSERT(x)) = "\assert{" x "}"
%format (ASSERT'(y)(x)) = "\varassert{" y "}{" x "}"
%format (ASSERTRANGE(x)) = "\assertrange{" x "}"
%format (ASSERTRANGE'(y)(x)) = "\varassertrange{" y "}{" x "}"

\newcommand{\awa}[2]{\mathrlap{#2}\phantom{#1}} % as wide as

%% Some recommended packages.
\usepackage{booktabs}   %% For formal tables:
                        %% http://ctan.org/pkg/booktabs
\usepackage{subcaption} %% For complex figures with subfigures/subcaptions
                        %% http://ctan.org/pkg/subcaption

\usepackage{mathtools}

\usepackage[color=yellow,textsize=footnotesize]{todonotes}
\setlength{\marginparwidth}{2cm}
\reversemarginpar

\renewcommand{\sectionautorefname}{Section}
\renewcommand{\subsectionautorefname}{Section}
\renewcommand{\subsubsectionautorefname}{Section}
\renewcommand{\figureautorefname}{Figure}

\makeatletter\if@@ACM@@journal\makeatother
%% Journal information (used by PACMPL format)
%% Supplied to authors by publisher for camera-ready submission
\acmJournal{PACMPL}
%\acmVolume{1}
%\acmNumber{1}
%\acmArticle{1}
\acmYear{2017}
\acmMonth{2}
%\acmDOI{10.1145/nnnnnnn.nnnnnnn}
%\startPage{1}
%\else\makeatother
%% Conference information (used by SIGPLAN proceedings format)
%% Supplied to authors by publisher for camera-ready submission
%\acmConference[PL'17]{ACM SIGPLAN Conference on Programming Languages}{January 01--03, 2017}{New York, NY, USA}
%\acmYear{2017}
%\acmISBN{978-x-xxxx-xxxx-x/YY/MM}
%\acmDOI{10.1145/nnnnnnn.nnnnnnn}
%\startPage{1}
\fi


%% Copyright information
%% Supplied to authors (based on authors' rights management selection;
%% see authors.acm.org) by publisher for camera-ready submission
\setcopyright{none}             %% For review submission
%\setcopyright{acmcopyright}
%\setcopyright{acmlicensed}
%\setcopyright{rightsretained}
%\copyrightyear{2017}           %% If different from \acmYear


%% Bibliography style
\bibliographystyle{ACM-Reference-Format}
%% Citation style
%% Note: author/year citations are required for papers published as an
%% issue of PACMPL.
\citestyle{acmauthoryear}   %% For author/year citations



\begin{document}

\setlength{\mathindent}{\parindent}

%% Title information
\title{An Axiomatic Basis for Bidirectional Programming}
                                        %% [Short Title] is optional;
                                        %% when present, will be used in
                                        %% header instead of Full Title.
%\titlenote{with title note}            %% \titlenote is optional;
                                        %% can be repeated if necessary;
                                        %% contents suppressed with 'anonymous'
%\subtitle{Subtitle}                    %% \subtitle is optional
%\subtitlenote{with subtitle note}      %% \subtitlenote is optional;
                                        %% can be repeated if necessary;
                                        %% contents suppressed with 'anonymous'


%% Author information
%% Contents and number of authors suppressed with 'anonymous'.
%% Each author should be introduced by \author, followed by
%% \authornote (optional), \orcid (optional), \affiliation, and
%% \email.
%% An author may have multiple affiliations and/or emails; repeat the
%% appropriate command.
%% Many elements are not rendered, but should be provided for metadata
%% extraction tools.

%% Author with single affiliation.
\author{Hsiang-Shang Ko}
%\authornote{with author1 note}         %% \authornote is optional;
                                        %% can be repeated if necessary
%\orcid{nnnn-nnnn-nnnn-nnnn}            %% \orcid is optional
\affiliation{
  \position{Project Researcher}
  \department{Information Systems Architecture Science Research Division}
                                        %% \department is recommended
  \institution{National Institute of Informatics}
                                        %% \institution is required
  \streetaddress{2-1-2 Hitotsubashi}
  \city{Chiyoda-ku}
  \state{Tokyo}
  \postcode{101-8430}
  \country{Japan}
}
\email{hsiang-shang@@nii.ac.jp}         %% \email is recommended

%% Author with two affiliations and emails.
\author{Zhenjiang Hu}
%\authornote{with author2 note}         %% \authornote is optional;
                                        %% can be repeated if necessary
%\orcid{nnnn-nnnn-nnnn-nnnn}            %% \orcid is optional
\affiliation{
  \position{Professor}
  \department{Information Systems Architecture Science Research Division}
                                        %% \department is recommended
  \institution{National Institute of Informatics}
                                        %% \institution is required
  \streetaddress{2-1-2 Hitotsubashi}
  \city{Chiyoda-ku}
  \state{Tokyo}
  \postcode{101-8430}
  \country{Japan}
}
\email{hu@@nii.ac.jp}         %% \email is recommended


%% Paper note
%% The \thanks command may be used to create a "paper note" ---
%% similar to a title note or an author note, but not explicitly
%% associated with a particular element.  It will appear immediately
%% above the permission/copyright statement.
%\thanks{with paper note}               %% \thanks is optional
                                        %% can be repeated if necesary
                                        %% contents suppressed with 'anonymous'


%% Abstract
%% Note: \begin{abstract}...\end{abstract} environment must come
%% before \maketitle command
\begin{abstract}
Text of abstract \ldots.
\end{abstract}


%% 2012 ACM Computing Classification System (CSS) concepts
%% Generate at 'http://dl.acm.org/ccs/ccs.cfm'.
%\begin{CCSXML}
%<ccs2012>
%<concept>
%<concept_id>10011007.10011006.10011008</concept_id>
%<concept_desc>Software and its engineering~General programming languages</concept_desc>
%<concept_significance>500</concept_significance>
%</concept>
%<concept>
%<concept_id>10003456.10003457.10003521.10003525</concept_id>
%<concept_desc>Social and professional topics~History of programming languages</concept_desc>
%<concept_significance>300</concept_significance>
%</concept>
%</ccs2012>
%\end{CCSXML}

%\ccsdesc[500]{Software and its engineering~General programming languages}
%\ccsdesc[300]{Social and professional topics~History of programming languages}
%% End of generated code


%% Keywords
%% comma separated list
%\keywords{keyword1, keyword2, keyword3}  %% \keywords is optional


%% \maketitle
%% Note: \maketitle command must come after title commands, author
%% commands, abstract environment, Computing Classification System
%% environment and commands, and keywords command.
\maketitle


\section{Introduction}
\label{sec:introduction}

\citet{Ko-BiGUL}

\citet{Foster-lenses}

\begin{code}
putWidth  : (Int, Int) <~ Int
putWidth  =
  rearrV v -> (v, ())
       replace
    *  skip const ()
\end{code}

\begin{code}
putSquare  : (Int, Int) <~ Int
putSquare  =
  case
    normal (w, _) v | w == v exit _
      skip fst
    adaptive _ _
      \ _ v -> (v, v)
\end{code}

\section{Theory of putback triples}
\label{sec:putback-triples}

\begin{theorem}[soundness of putback triples] \label{thm:putback-soundness}
\[ |(ASSERT R) ~ b ~ (ASSERT R')| \quad|=>|\quad |(all(s, v)) R s v => (some(s')) put b s v = Just s' && R' s' s v| \]
\end{theorem}

\begin{theorem}[partial forward consistency]
\[ |(ASSERT R) ~ b ~ (ASSERT(s' _ v || R' s' v))| \quad|=>|\quad |(all(s, v)) get b s = Just v && R s v ~ => R' s v| \]
\end{theorem}
\begin{proof}
Suppose |get b s = Just v| and |R s v|.
By GetPut, we have |put b s v = Just s|, and by \autoref{thm:putback-soundness}, we have |put b s v = Just s' && R' s' v|.
Thus |s = s'|, and we obtain |R' s v|.
\end{proof}

\section{Putback proof rules}
\label{sec:putback-proof-rules}

\[ \begin{array}{l} \hline
|(ASSERT empty) ~ fail ~ (ASSERT empty)|
\end{array} \]

\[ \begin{array}{l} \hline
|(ASSERT(_ _)) ~ replace ~ (ASSERT(s' _ v || s' = v))|
\end{array} \]

\[ \begin{array}{l} \hline
|(ASSERT(s v || f s = v)) ~ skip f ~ (ASSERT(s' s _ || s' = s))|
\end{array} \]

\[ \begin{array}{l}
|{-"\assert{\awa{T}{R}}\;"-} ~ {-"\awa{r}{l}\;\;"-} ~ {-"\assert{\awa{T^\prime}{R^\prime}}"-}| \\
|{-"\assert{T}\;"-} ~ {-"r\;\;"-} ~ {-"\assert{T^\prime}"-}| \\ \hline
|(ASSERT(R * T)) ~ l * r ~ (ASSERT(R' * T'))|
\end{array} \]

can indeed be regarded as a simpler variant of separating conjunction~\citep{Reynolds-separation-logic}

\[ \begin{array}{l}
|(ASSERT(s wpat || R s (ENV(wpat)))) ~ b ~ (ASSERT(s' s wpat || R' s' s (ENV(wpat))))| \\ \hline
|(ASSERT(s vpat || R s (ENV(vpat)))) ~ rearrV vpat -> wpat; ~ b ~ (ASSERT(s' s vpat || R' s' s (ENV(vpat))))|
\end{array} \]

\[ \begin{array}{l}
|(ASSERT(tpat v || R (ENV(tpat)) v)) ~ b ~ (ASSERT(tpat' tpat v || R' (ENV(tpat')) (ENV(tpat)) v))| \\ \hline
|(ASSERT(spat v || R (ENV(spat)) v)) ~ rearrS spat -> tpat; ~ b ~ (ASSERT(spat' spat v || R' (ENV(spat')) (ENV(spat)) v))|
\end{array} \]

\[ \begin{array}{l}
|(ASSERT R) ~ b ~ (ASSERT R')| \\
|T => R| \\
|all(s', s, v) R' s' s v && T s v => T' s' s v| \\ \hline
|(ASSERT T) ~ b ~ (ASSERT T')|
\end{array} \]

\begin{code}
(ASSERT(_ _))
rearrV v -> (v, ())
  (ASSERT(_ (_ , ())))
     (ASSERT(_ _))
     replace
     (ASSERT(w' _ v | w' = v)) (ASSERTRANGE(_))
  *  (ASSERT(_ ()))
     (ASSERT(h () | const () h = ()))
     skip const ()
     (ASSERT(h' h () | h' = h)) (ASSERTRANGE(_))
  (ASSERT((w', h') (_ , h) (v, ()) | w' = v && h' = h)) (ASSERTRANGE(_))
(ASSERT((w', h') (_ , h) v | w' = v && h' = h)) (ASSERTRANGE(_))
\end{code}

\begin{code}
(ASSERT(_ _))
case
  normal (w, _) v | w == v exit _
    (ASSERT((w, _) v | w = v))
    (ASSERT((w, h) v | fst (w, h) = v))
    skip fst
    (ASSERT((w', h') (w, h) _ | (w', h') = (w, h))) (ASSERTRANGE(_))
    (ASSERT((w', h') (w, h) v | w' = v && (w = v => h' = h) && (w /= v => w' = h')))
  adaptive _ _
    (ASSERT(not (COM((w, _) v | w = v))))
    (ASSERT((w, _) v | w /= v))
    \ (w, h) v -> (v, v)
    (ASSERT'(reentry)(COM((w, _) v | w = v) (v, v) v))
    (ASSERT'(retentiveness)(all(w', h') Post (w', h') (v, v) v => Post (w', h') (w, h) v))
  (ASSERT'(coverage)(COM((w, _) v | w = v) || (COM(_ _)) ))
  (ASSERT'(disjointness)(True))
(ASSERT(Post @ (COM((w', h') (w, h) v | w' = v && (w = v => h' = h) && (w /= v => w' = h'))))) (ASSERTRANGE(_))
\end{code}

\begin{code}
replaceHead =
  (ASSERT((_ :: _) _))
  rearrS (s::ss) -> (s, ss)
    (ASSERT((_ , _) _))
    rearrV v -> (v, ())
      (ASSERT((_ , _) (_ , ())))
         (ASSERT(_ _))
         replace
         (ASSERT(s' _ v | s' = v)) (ASSERTRANGE(_))
      *  (ASSERT(_ ()))
         (ASSERT(ss () | const () ss = ()))
         skip const ()
         (ASSERT(ss' ss _  | ss' = ss)) (ASSERTRANGE(_))
         (ASSERT(ss' ss () | ss' = ss))
      (ASSERT((s', ss') (s, ss) (v, ()) | s' = v && ss' = ss)) (ASSERTRANGE((_ , _)))
    (ASSERT((s', ss') (s, ss) v | s' = v && ss' = ss)) (ASSERTRANGE((_ , _)))
  (ASSERT((s'::ss') (s::ss) v | s' = v && ss' = ss)) (ASSERTRANGE((_ :: _)))
  (ASSERT((s'::ss') ss v | s' = v && ss' = tail ss))
\end{code}

\begin{code}
(ASSERT(_ _))
case
  normal (_ :: _) _ exit _
    (ASSERT((_ :: _) _))
    replaceHead
    (ASSERT((s'::ss') ss v | s' = v && ss' = tail ss)) (ASSERTRANGE((_ :: _)))
  adaptive _ _
    (ASSERT(not (COM((_ :: _) _))))
    (ASSERT([] _))
    \ _ v -> [v]
    (ASSERT'(reentry)(COM((_ :: _) _) [v] v))
    (ASSERT'(retentiveness)(all(ss') Post ss' [v] v => Post ss' [] v))
  (ASSERT'(coverage)(COM((_ :: _) _) || (COM(_ _))))
  (ASSERT'(disjointness)(True))
(ASSERT(Post @ (COM((s'::ss') ss v | s' = v && ss' = tail ss)))) (ASSERTRANGE((_ :: _)))
\end{code}

\section{Theory of range triples}
\label{sec:range-triples}

\begin{theorem}[soundness of range triples]
\[ |(ASSERTRANGE R) ~ b ~ (ASSERTRANGE P)| \quad|=>|\quad |(all(s)) P s => (some(v)) get b s = Just v && R s v| \]
\end{theorem}
The conclusion implies (by GetPut)
\[ |(all(s)) P s => (some(v)) put b s v = Just s && R s v| \]
which further implies
\[ |(all(s')) P s' => (some(s, v)) put b s v = Just s' && R s v| \]

\begin{theorem}[total forward consistency] \label{thm:total-forward-consistency}
\[ |(ASSERT R) ~ b ~ (ASSERT(s' _ v || R' s' v))| \quad\wedge\quad |(ASSERTRANGE R) ~ b ~ (ASSERTRANGE P)| \quad|=>|\quad |(all(s)) P s => (some(v)) get b s = Just v && R' s v| \]
\end{theorem}

\section{Range proof rules}
\label{sec:range-proof-rules}

\[ \begin{array}{l}
|(ASSERTRANGE R) ~ b ~ (ASSERTRANGE P)| \\
|all(s, v) R s v && Q s => T s v| \\
|Q => P| \\ \hline
|(ASSERTRANGE T) ~ b ~ (ASSERTRANGE Q)|
\end{array} \]

\[ \begin{array}{l}
\hline
|(ASSERTRANGE empty) ~ fail ~ (ASSERTRANGE empty)|
\end{array} \]

\[ \begin{array}{l}
\hline
|(ASSERTRANGE(s v || f s = v)) ~ skip f ~ (ASSERTRANGE _)|
\end{array} \]

\[ \begin{array}{l}
\hline
|(ASSERTRANGE(s v || s = v)) ~ replace ~ (ASSERTRANGE _)|
\end{array} \]

\[ \begin{array}{l}
|{-"\assertrange{\awa{T}{R}}\;"-} ~ {-"\awa{r}{l}\;\;"-} ~ {-"\assertrange{\awa{Q}{P}}"-}| \\
|{-"\assertrange{T}\;"-} ~ {-"r\;\;"-} ~ {-"\assertrange{Q}"-}| \\ \hline
|(ASSERTRANGE(R * T)) ~ l * r ~ (ASSERTRANGE(P * Q))|
\end{array} \]

\[ \begin{array}{l}
|(ASSERTRANGE(s wpat || R s (ENV(wpat)))) ~ b ~ (ASSERTRANGE P)| \\ \hline
|(ASSERTRANGE(s vpat || R s (ENV(vpat)))) ~ rearrV vpat -> wpat; ~ b ~ (ASSERTRANGE P)|
\end{array} \]

\[ \begin{array}{l}
|(ASSERTRANGE(tpat v || R (ENV(tpat)) v)) ~ b ~ (ASSERTRANGE(tpat || P (ENV(tpat))))| \\ \hline
|(ASSERTRANGE(spat v || R (ENV(spat)) v)) ~ rearrS spat -> tpat; ~ b ~ (ASSERTRANGE(spat || P (ENV(spat))))|
\end{array} \]

\begin{lemma}
\[ |(ASSERT R) ~ b ~ (ASSERT(s' _ _ || Q s')) (ASSERTRANGE P)| \quad|=>|\quad |P => Q| \]
\end{lemma}

\begin{code}
putSquare' =
  (ASSERT(_ _))
  rearrV v -> (v, v)
    (ASSERT(_ (v, v)))
    (ASSERT'(weaken)(_ _))
       (ASSERT(_ _))
       replace
       (ASSERT(x' _ v | x' = v))
    *  (ASSERT(_ _))
       replace
       (ASSERT(y' _ w | y' = w))
    (ASSERT((x', y') _ (v, w) | x' = v && y' = w))
    (ASSERT((x', y') _ (v, v) | x' = v && y' = v))
  (ASSERT((x', y') _ v | x' = v && y' = v))
\end{code}

\begin{code}
(ASSERTRANGE(_ (v, w) | v = w))

(ASSERTRANGE((s, t) (v, w) | s = v && t = w))
   (ASSERTRANGE(s v | s = v))
   replace
   (ASSERTRANGE _)
*  (ASSERTRANGE(t w | t = w))
   replace
   (ASSERTRANGE _)
(ASSERTRANGE _)
(ASSERTRANGE((s, t) | s = t))
\end{code}

\section{Recursion}
\label{sec:recursion}

$b \preceq c$ exactly when |(all(s, v, s')) put b s v = Just s' => put c s v = Just s'|

|get b s = Just v ~ => ~ put b s v = Just s ~ => ~ put c s v = Just s ~ => ~ get c s = Just v|

$( b_i \mid i \in \mathbb N )$ is an increasing chain exactly when $i \leq j \Rightarrow b_i \preceq b_j$

|put ((LUB(i)) b_i) s v = Just s'| exactly when |put b_i s v = Just s'| for some~|i|

A function~$f$ is \emph{monotonic} exactly when $b \preceq c$ implies $|f b| \preceq |f c|$

A function~$f$ is \emph{continuous} exactly when |(LUB(i)) (f b_i) = f((LUB(i)) b_i)|

|put ((LUB(i)) (f b_i)) s v = Just s'| if and only if |put (f((LUB(i)) b_i)) s v = Just s'|

|put (f b_i) s v = Just s'| for some~|i| if and only if |put (f((LUB(i)) b_i)) s v = Just s'|

\begin{code}
replaceAll : Eq A => ~ [A] <~ A
replaceAll =
  case
    adaptive [] x
      \ _ x -> [x]
    normal (_ :: []) _ exit (_ :: [])
      rearrS (s :: []) -> s
        replace
    normal (_ :: _) _ exit (_ :: _ :: _)
      rearrS (s :: ss) -> (s, ss)
        rearrV v -> (v, v)
             replace
          *  replaceAll
\end{code}

\begin{code}
(ASSERT(_ _))
replaceAll
(ASSERT(ss' ss v | ((all(s')) s' `elem` ss' => s' = v) && (ss /= [] => length ss' = length ss)))
(ASSERTRANGE(ss | (all(x, y)) x `elem` ss && y `elem` ss => x = y))
\end{code}

Note that consistency is specified purely relationally in this case.
This reinforces the statement that it is the consistency relation that we have in mind before programming |put|, rather than the |get| function: for asymmetric lenses, a consistency relation must be functional (in the sense that every source is related to at most one view) but may well not be executable (as in the case of |replaceAll|), whereas a |get| function is executable, and its computational content is non-trivially derived from the |put| program.
\todo{Discuss asymmetric lenses from \citet{Stevens-QVT}'s perspective somewhere before, or move the entire remark to \autoref{sec:discussion}.}
\begin{proposition}
\[ |(ASSERT R) ~ b ~ (ASSERT(s' s _ || s' = s))| \quad|=>|\quad |(all(s, v)) R s v => get b s = Just v| \]
\end{proposition}

\section{Verifying key-based alignment}
\label{sec:key-based-alignment}

Several variants of alignment have been implemented in BiGUL~\citep{Zan-Brul, Mendes-delta-alignment}.
Here we choose to verify a variant that is non-trivial and yet not overly complicated: key-based alignment.

For now, \citet{Barbosa-matching-lenses}'s matching lenses may be more powerful, but they are special-purpose and require the invention of a new framework in which to build everything from scratch.
Instead, we strive to implement complex lenses from just a fixed set of primitives, treating bidirectional programming as having the same status of general-purpose programming.

\begin{code}
keyAlign : Eq K => (S -> K) -> (V -> K) -> (S <~ V) -> (V -> S) -> ([S] <~ [V])
keyAlign ks kv b c =
  case
    normal [] [] exit []
      rearrV [] -> ()
         skip const ()
    normal (s :: _) (v :: _) | ks s == kv v exit (_ :: _)
      rearrS (s :: ss) -> (s, ss)
        rearrV (v :: vs) -> (v, vs)
             b
          *  keyAlign ks kv b c
    adaptive (_ :: _) []
      \ _ _ -> []
    adaptive ss (v :: vs) | kv v `elem` map ks ss
      \ ss (v :: _) -> extract ks kv v ss
    adaptive _ (_ :: _)
      \ ss (v :: _) -> c v :: ss
\end{code}

\begin{code}
extract : Eq K => (S -> K) -> (V -> K) -> V -> [S] -> [S]
extract ks kv v []         = []
extract ks kv v (s :: ss)  | ks s == kv v  =  s :: ss
                           | otherwise     =  case extract ks kv v ss of
                                                []         -> s :: []
                                                s' :: ss'  -> s' :: s :: ss'
\end{code}

\section{Discussion}
\label{sec:discussion}

A proper characterisation of the behaviour of composition will be essential for avoiding problems about composition as discussed by, e.g., \citet{Diskin-delta-asymmetric}.
\[ \begin{array}{l}
|(ASSERT (a b' || (some(b, c)) R a c && R' a b && U' b' b c)) ~ l ~ (ASSERT(COM(a' _ b || R' a' b) && T') (ASSERTRANGE P)| \\
|(ASSERT (b c || (some(a)) R' a b && R a c)) ~ r ~ (ASSERT U')|
\\ \hline
|(ASSERT(R && (COM(a _ || P a)))) ~ l . r ~ (ASSERT(a' a c || (some(b, b')) T' a' a b' && U' b' b c ))|
\end{array} \]

We need general and formal reasoning principles for bidirectional programming to make it become a proper discipline of programming.
Within the realm of asymmetric lenses, the get-based approach is fundamentally just about \emph{programming consistency specifications}; only with the putback-based approach can we start talking about \emph{programming to meet functional specifications}.

Programming languages should be shipped with reasoning principles; even domain-specific languages deserve domain-specific reasoning principles!

Doesn't \autoref{thm:total-forward-consistency} say that we need to reason in both directions anyway?
Even if the sceptic reader insists that constructing the range triples requires thinking in the forward direction, they will still have to admit that, according to \autoref{thm:total-forward-consistency}, what needs to be verified in the forward direction is normally much easier than proving that |get| is contained in the consistency relation. This is because the major work has been done in the backward direction.

It has been argued that ``putback'' is the essence of bidirectional programming~\citep{Fischer-essence-of-BX}.
We would like to amend this statement: Putback-based \emph{reasoning} is the essence of bidirectional programming.
This paper has shown, with the Hoare-style logic, what that reasoning can look like.

% Acknowledgements
\begin{acks}                            %% acks environment is optional
                                        %% contents suppressed with 'anonymous'
  %% Commands \grantsponsor{<sponsorID>}{<name>}{<url>} and
  %% \grantnum[<url>]{<sponsorID>}{<number>} should be used to
  %% acknowledge financial support and will be used by metadata
  %% extraction tools.
  Jeremy Gibbons and Akimasa Morihata
  
  This material is based upon work supported by the
  \grantsponsor{GS100000001}{National Science
    Foundation}{http://dx.doi.org/10.13039/100000001} under Grant
  No.~\grantnum{GS100000001}{nnnnnnn} and Grant
  No.~\grantnum{GS100000001}{mmmmmmm}.
\end{acks}

%% Bibliography
\bibliography{../../../Documents/bib}


%% Appendix
%\appendix
%\section{Appendix}
%
%Text of appendix \ldots

\end{document}
